\documentclass[11pt, a4paper]{amsart}

\usepackage{amsmath, amssymb, amsthm}
\usepackage{enumitem}
\usepackage{hyperref}
\usepackage{booktabs}
\usepackage{array}
\usepackage[dvipsnames]{xcolor}

\theoremstyle{plain}
\newtheorem{theorem}{Theorem}[section]
\newtheorem{proposition}[theorem]{Proposition}
\newtheorem{lemma}[theorem]{Lemma}
\newtheorem{corollary}[theorem]{Corollary}

\theoremstyle{definition}
\newtheorem{definition}[theorem]{Definition}
\newtheorem{remark}[theorem]{Remark}
\newtheorem{hypothesis}[theorem]{Hypothesis}

\newcommand{\MI}{\operatorname{MI}}
\newcommand{\fS}{\mathfrak{S}}
\newcommand{\cL}{\mathcal{L}}
\newcommand{\eps}{\varepsilon}

\title[A Burgers-Fisher Equation for Prime Gaps]{%
A Burgers-Fisher Equation for Prime Gaps:\\
Spectral Analysis of Convergence to Poisson}

\author{Tobias Canavesi}
\address{Independent Researcher}
\email{tobiascanavesi@gmail.com}

\date{April 2026}

\begin{document}

\begin{abstract}
The emergence of prime numbers from the deterministic sieve of Eratosthenes 
into random probability at the macroscopic level, described by a Poisson process, 
represents one of the most basic yet most interesting phenomena in number theory. 
In this paper, we develop a new physical mechanism to describe this transition based on modeling the time-evolving 
distribution of prime gaps as fluids moving through a porous medium. If the normalized 
gap density follows a Feller-type Fokker-Planck equation---chosen to reflect the strict 
constraints imposed by the nature of primes on short intervals---we show that the 
corresponding velocity field satisfies a nonlinear Burgers-Fisher PDE. For our proposed 
system, the Poisson regime does not represent an arbitrary statistical limit; 
it is the only stable fixed point of the PDE, and it is stabilized by a Fisher-KPP 
reaction term. By performing a spectral analysis of the linearized operator, we find 
that there exists a Laguerre basis which forms an eigenbasis for the operator, 
and contains a spectral gap of size $1$, indicating that any perturbation forced 
by a decaying function will be removed to order $O(1 / (\log x)^\beta)$. Lastly, we have 
conducted numerical simulations across eleven scales up to $x=10^{8}$, and demonstrated 
that our theoretical framework provides a good description of both the decay rate of 
the mutual information between successive prime gaps as $(\log x)^{-0.19}$, and the fact 
that nearly all of the information content in this measure can be attributed to mod-$3$ 
sieve constraints (approximately $94\%$).
\end{abstract}

\maketitle
\tableofcontents

% ====================================================================
\section{Introduction}
% ====================================================================

We know the set of all prime numbers follows strictly deterministic laws and is defined completely by the sieve of Eratosthenes. 
However, at large scales, the gaps between subsequent prime numbers act exactly like a purely memoryless random process: Gallagher~\cite{Gallagher1976} 
has shown, conditionally on the Hardy-Littlewood conjecture~\cite{HL1923}, that the counts of primes in short intervals will be Poisson distributed 
and therefore the normalized gaps should follow an exponential distribution $f(v)=e^{-v}$. How does a system based solely on rigid divisibility laws 
become statistically random?

Although the \emph{end point} (Poisson) is conditionally known, the \emph{process} which leads to this convergence remains unknown. 
Lemke Oliver and Soundararajan~\cite{LO2016} showed that there is a statistically significant bias for the residue classes of consecutive primes, 
indicative of a non-negligible residual influence of the sieve, decaying as $O((\log \log x)/(\log x))$; Holt~\cite{Holt2024} attempted to 
model these biases using sieve dynamics. However, to date, no differential equation modeling the time evolution or quantifying the rate of 
convergence has been proposed.

In this work, we believe that the explanation for this behavior can be found in the language of fluid mechanics. Instead of viewing the density of 
prime gaps as a static probability distribution, we represent it as a \emph{velocity field} and view our system as a fluid flowing through a porous 
medium---the sieve. We demonstrate that if the density of the gaps evolves according to a Feller Fokker-Planck equation---whose state-dependent diffusion 
represents the arithmetical rigidity of small gaps---then its corresponding velocity field satisfies a nonlinear Burgers-Fisher 
equation~\cite{Fisher1937,KPP1937}.

The Poisson regime is not simply a statistical accident; rather, it is the one stable solution to this partial differential equation; 
this stability is guaranteed by both a spectral gap of $1$ and a Fisher-KPP reaction term $u(1-u)$, generated naturally by the structure of 
the Fokker-Planck equation. We verified numerically the convergence at $11$ different length scales with a decay rate $\approx (\log x)^{-0.19}$; 
however, it was clear that the dominant contribution came primarily from the single prime $p=3$.

% ====================================================================
\section{The Burgers--Fisher Equation for Prime Gaps}
\label{sec:equation}
% ====================================================================

\subsection{The Feller process and its Fokker--Planck equation}

We want to find an appropriate diffusion process whose stationary distribution is given by $\mathrm{Exp}(1)$, 
the so-called Gallagher-Poisson distribution. The Feller diffusion process is, among all the diffusion processes with linear drift, 
the only one where we obtain both the required form for the diffusion coefficient in order to reflect the empirical structure of prime gaps, 
as well as the correct behavior of the fluctuations. Indeed, since the noise of the Feller diffusion process is given by $\sigma(V) = \sqrt{2V}$, 
it follows that small gaps ($V \approx 0$) will be almost deterministic, whereas large gaps will behave like free random variables. 
This reflects the strong arithmetic constraints imposed by the sieve on the small gap values (i.e., only $g=2$, $g=4$, and $g=6$ are common), 
whereas large gap values can follow a much broader distribution. Cohen~\cite{Cohen2024} independently confirmed that the moments of prime gaps 
converge to those of the exponential distribution, further supporting this choice.

Consider the Feller diffusion (Cox--Ingersoll--Ross process)
\begin{equation}\label{eq:sde}
dV = (1 - V)\,dt + \sqrt{2V}\,dW_t, \qquad V \ge 0,
\end{equation}
with mean-reversion toward~$1$ and state-dependent noise $\sigma(V) = \sqrt{2V}$.
Its Fokker--Planck equation for the density $f(v,t)$ is
\begin{equation}\label{eq:fp}
\partial_t f = v f'' + (v+1)f' + f,
\end{equation}
where primes denote $\partial_v$.

\begin{lemma}\label{lem:stationary}
$f_\infty(v) = e^{-v}$ is the unique probability density satisfying $\cL f_\infty = 0$, where $\cL f = vf'' + (v+1)f' + f$.
\end{lemma}

\begin{proof}
Direct verification: $\cL(e^{-v}) = v \cdot e^{-v} + (v+1)(-e^{-v}) + e^{-v} = 0$.
Uniqueness follows from the general theory of Feller diffusions~\cite{Risken1989}: 
the process~\eqref{eq:sde} with $2 \times \text{mean-reversion speed} / \sigma^2 = 2 \cdot 1 / 2 = 1 \ge 1$ 
has an entrance boundary at $0$ and a natural boundary at $\infty$, 
admitting a unique stationary distribution.
\end{proof}

\subsection{Derivation of the velocity equation}

\begin{theorem}[The sieve Burgers--Fisher equation]
\label{thm:burgers-fisher}
Let $f(v,t) > 0$ satisfy the Fokker--Planck equation~\eqref{eq:fp} on $(0,\infty)$.
Define the velocity field
\begin{equation}\label{eq:velocity}
u(v,t) = -\frac{\partial}{\partial v}\log f(v,t) = -\frac{f'}{f}.
\end{equation}
Then $u$ satisfies the Burgers--Fisher equation in a porous medium:
\begin{equation}\label{eq:burgers-fisher}
\boxed{\partial_t u + 2v\,u\,\partial_v u = v\,\partial_v^2 u + (v{+}2)\,\partial_v u + u(1{-}u).}
\end{equation}
\end{theorem}

\begin{proof}
Set $\psi = \log f$, so $\psi' = f'/f$ and $f''/f = \psi'' + (\psi')^2$.
From~\eqref{eq:fp}:
\begin{equation}\label{eq:psi}
\psi_t = v(\psi'' + \psi'^2) + (v+1)\psi' + 1.
\end{equation}
Differentiate~\eqref{eq:psi} with respect to $v$:
\begin{align*}
\psi'_t &= \psi'' + v\psi''' + \psi'^2 + 2v\psi'\psi'' + \psi' + (v+1)\psi'' \\
&= v\psi''' + (v+2)\psi'' + 2v\psi'\psi'' + \psi'^2 + \psi'.
\end{align*}
Substitute $u = -\psi'$ (hence $\psi'' = -u'$, $\psi''' = -u''$):
\begin{align*}
-u_t &= -vu'' - (v+2)u' + 2vu\,u' + u^2 - u.
\end{align*}
Multiply by $-1$ and rearrange:
\[
u_t = vu'' + (v+2)u' - 2vu\,u' - u^2 + u = vu'' + (v+2)u' + u(1-u) - 2vu\,u'. \qedhere
\]
\end{proof}

\begin{remark}[Structure of the equation]
Equation~\eqref{eq:burgers-fisher} combines three classical PDE mechanisms:
\begin{enumerate}
\item \emph{Burgers nonlinearity} $2vu\,u_v$: the hallmark of the Navier--Stokes and Burgers equations, 
arising here from the nonlinear transformation $u = -(\log f)'$.
\item \emph{Porous medium diffusion} $v\,u_{vv}$: the diffusivity equals $v$, so small gaps ($v \approx 0$) are ``tight'' 
and large gaps fluctuate freely.
\item \emph{Fisher--KPP reaction} $u(1{-}u)$: this term pushes $u$ toward~$1$ from either side. It emerges from the Fokker--Planck structure, not imposed by hand.
\end{enumerate}
\end{remark}


% ====================================================================
\section{Spectral Stability of the Poisson Equilibrium}
\label{sec:spectral}
% ====================================================================

\begin{theorem}[Spectral analysis]
\label{thm:spectral}
\leavevmode
\begin{enumerate}[label=\upshape(\roman*)]
\item The unique bounded steady state of~\eqref{eq:burgers-fisher} is $u \equiv 1$, corresponding to the 
exponential distribution $f = e^{-v}$.
\item The linearization around $u = 1$, writing $\delta u = u - 1$, is
\begin{equation}\label{eq:linearized}
\partial_t(\delta u) = v\,(\delta u)'' + (2 - v)\,(\delta u)' - \delta u.
\end{equation}
\item The eigenvalue problem $v\,h'' + (2-v)\,h' + (\lambda+1)\,h = 0$ is the associated Laguerre 
equation with parameter $\alpha = 1$.
Its solutions are the associated Laguerre polynomials $L_n^{(1)}(v)$ with eigenvalues
\begin{equation}\label{eq:eigenvalues}
\lambda_n = -(n+1), \qquad n = 0, 1, 2, \ldots
\end{equation}
\item The spectral gap is $|\lambda_0| = 1$.
\end{enumerate}
\end{theorem}

\begin{proof}
(i) A constant steady state $u = c$ satisfies $c(1-c) = 0$, giving $c = 0$ or $c = 1$.
If $u = 0$ then $f' = 0$, i.e., $f = \text{const}$, which is not integrable on $[0,\infty)$.
Thus $u = 1$ is the unique physically admissible steady state.

(ii) Substitute $u = 1 + \delta u$ into~\eqref{eq:burgers-fisher}, keeping linear terms in $\delta u$:
\begin{align*}
(\delta u)_t + 2v \cdot 1 \cdot (\delta u)' &= v(\delta u)'' + (v+2)(\delta u)' + (1 + \delta u)(1 - 1 - \delta u) \\
&= v(\delta u)'' + (v+2)(\delta u)' - \delta u + O((\delta u)^2).
\end{align*}
Hence $(\delta u)_t = v(\delta u)'' + (2-v)(\delta u)' - \delta u$.

(iii) Setting $(\delta u)(v,t) = h(v)\,e^{\lambda t}$ gives $v h'' + (2-v)h' + (\lambda + 1)h = 0$.
With the substitution $\alpha = 1$ and $n = \lambda + 1$, this is the associated Laguerre equation $v h'' + (\alpha + 1 - v)h' + nh = 0$, whose polynomial solutions are $L_n^{(1)}(v)$~\cite{Szego1975}.

(iv) From $\lambda_n = -(n+1)$: $\lambda_0 = -1$, so all perturbation modes decay at rate $\ge 1$.
\end{proof}


% ====================================================================
\section{Convergence Under Decaying Sieve Forcing}
\label{sec:convergence}
% ====================================================================

\begin{theorem}[Perturbed convergence]
\label{thm:convergence}
Let $u(v,t)$ satisfy
\begin{equation}\label{eq:forced}
\partial_t u + 2vu\,u_v = vu_{vv} + (v+2)u_v + u(1-u) + \eps(t)\,F(v),
\end{equation}
where $\|F\|_\infty \le M$ and $|\eps(t)| \le C_0/t^\beta$ for $\beta > 0$.
Then there exists $C > 0$ such that $\|\delta u(\cdot, t)\|_{L^2(v\,e^{-v})} \le C/t^\beta$ for all $t$ sufficiently large.
\end{theorem}

\begin{proof}
Define the energy $E(t) = \frac{1}{2}\|\delta u\|^2_{L^2(v e^{-v})}$, where $L^2(ve^{-v})$ is the natural weighted space for the associated Laguerre polynomials.

\textbf{Step 1 (Dissipation).} By the spectral gap (Theorem~\ref{thm:spectral}(iv)):
$\langle \cL_{\mathrm{lin}}(\delta u), \delta u \rangle \le -|\lambda_0| \cdot \|\delta u\|^2 = -\|\delta u\|^2.$

\textbf{Step 2 (Forcing).} $|\langle \eps F, \delta u \rangle| \le |\eps| \cdot M' \cdot \|\delta u\|$,
where $M' = \|F\|_{L^2(ve^{-v})}$.
By Young's inequality: $|\eps| M' \|\delta u\| \le \frac{1}{2}\|\delta u\|^2 + \frac{M'^2}{2}\eps^2$.

\textbf{Step 3 (Energy inequality).}
$\dot{E} \le -2E + E + \frac{M'^2}{2}\eps^2 = -E + \frac{C_1}{t^{2\beta}}.$

\textbf{Step 4 (Gr\"onwall).} The integrating factor $e^t$ gives
$E(t) \le E(t_0)e^{-(t-t_0)} + C_1\int_{t_0}^t e^{-(t-\tau)} \tau^{-2\beta}\,d\tau$.
The integral is $O(t^{-2\beta})$ by repeated integration by parts, yielding $E(t) = O(t^{-2\beta})$.
\end{proof}


% ====================================================================
\section{Connection to Prime Gaps}
\label{sec:primes}
% ====================================================================

\begin{hypothesis}[Sieve diffusion]
\label{hyp:sieve}
Under the Hardy--Littlewood conjecture, the normalized gap density $f(v, s)$ at scale $s = \log x$ 
can be written as $f = e^{-v}[1 + \varphi(v,s)]$, where the perturbation $\varphi$ is driven by the 
singular series and satisfies $\|\varphi(\cdot,s)\| = O(1/s^\beta)$ for some $\beta > 0$.
\end{hypothesis}

Under this hypothesis, the velocity field $u = 1 - \varphi'/(1+\varphi) \approx 1 - \varphi'$ satisfies a 
forced version of~\eqref{eq:burgers-fisher}, and Theorem~\ref{thm:convergence} gives $u \to 1$, i.e., 
convergence to Poisson.

\begin{remark}[What the hypothesis does not give]
Hypothesis~\ref{hyp:sieve} concerns the marginal gap density.
The mutual information $\MI(g_n; g_{n+1})$---a joint quantity---requires additional input.
We verify the MI decay numerically (\S\ref{sec:numerics}) but leave the rigorous joint implication as an 
open problem.
\end{remark}

\subsection{The sieve as porous medium}

Under the Hardy--Littlewood conjecture~\cite{HL1923}, the pair singular series gives:
$f(v,s) \propto \fS(\{0, vs\}) \cdot e^{-v}$,
where the relative weight is $W(g) = \fS(\{0,g\})/(2C_2) = \prod_{p|g,\,p>2} (p{-}1)/(p{-}2)$.
Each prime $p$ is a ``pore'' with strength $\delta_p = 1/(p{-}2)$: $\delta_3 = 1$, $\delta_5 = 1/3$, $\delta_7 = 1/5$.
As $s \to \infty$, gap values $g = vs$ grow and the mod-$p$ structure equidistributes, weakening each pore.

\subsection{Exact computation: the mod-$3$ coupling}

\begin{proposition}\label{prop:mod3}
Under i.i.d.\ uniform prime residues mod~$3$:
$\MI(g_n \bmod 3;\, g_{n+1} \bmod 3) = 1/4$ bit,
with structural zeros at $(1,1)$ and $(2,2)$ in the joint table.
\end{proposition}

\begin{proof}
With $r_n = p_n \bmod 3 \in \{1,2\}$ i.i.d.\ uniform, the gap $g_n \equiv r_{n+1} - r_n \pmod{3}$.
Enumerating all $8$ triples $(r_n, r_{n+1}, r_{n+2}) \in \{1,2\}^3$ yields the joint distribution:

\smallskip
\begin{center}
\begin{tabular}{c|ccc}
$g_n \bmod 3$ \textbackslash\ $g_{n+1} \bmod 3$ & 0 & 1 & 2 \\
\hline
0 & 2/8 & 1/8 & 1/8 \\
1 & 1/8 & \textbf{0} & 1/8 \\
2 & 1/8 & 1/8 & \textbf{0}
\end{tabular}
\end{center}

\smallskip\noindent
Marginals: $P(0) = 1/2$, $P(1) = P(2) = 1/4$. Only entries $(1,2)$ and $(2,1)$ deviate from the product 
of marginals:
$\MI = 2 \times \tfrac{1}{8}\log_2(\tfrac{1/8}{1/4 \cdot 1/4}) = \tfrac{1}{4}$ bit.
\end{proof}


% ====================================================================
\section{Numerical Verification}
\label{sec:numerics}
% ====================================================================

All measurements use $N = 200{,}000$ primes at 11 scales ($10^3$ to $10^8$).

\subsection{MI decay rate}
$\MI_{\mathrm{total}}(x) \sim 0.555 / (\log x)^{0.186}$, $R^2 = 0.986$.
To our knowledge, this is the first numerical measurement of the MI decay rate for prime gaps.

\subsection{Prime $3$ dominance}
$\MI_3 / \MI_{\mathrm{total}} \approx 94\%$ at all scales.
Empirical $\MI_3 = 0.302$ bits vs.\ theoretical $0.250$ bits; the $21\%$ excess over the i.i.d.\ prediction 
is consistent with the known biases in consecutive prime residues documented by Lemke Oliver and 
Soundararajan~\cite{LO2016}.
The structural zeros $P(1,1) = P(2,2) = 0$ hold \emph{exactly} for all primes $> 3$: if $g_n \equiv 1 \pmod{3}$ 
then $p_{n+1} \equiv 2 \pmod{3}$, forcing $g_{n+1} \equiv 1$ to require $p_{n+2} \equiv 0 \pmod{3}$, 
which is impossible for $p_{n+2} > 3$. Verified empirically with zero exceptions in $78{,}493$ 
consecutive gap pairs.

\subsection{Perturbation energy}
$E(s) = \|\varphi\|^2_{L^2(e^v)} \sim C/(\log x)^{0.63}$, $R^2 = 0.67$.
The Laguerre decomposition shows the $n=2$ mode (variance correction) dominates, while the $n=1$ 
mode (mean shift) is nearly zero.

\subsection{Information cascade}
\begin{center}
\begin{tabular}{cccc}
\toprule
Prime $p$ & Decay exponent $\alpha_p$ & $R^2$ & MI at $10^8$ \\
\midrule
3 & 0.22 & 0.980 & 0.302 bits \\
7 & 3.83 & 0.996 & 0.006 bits \\
11 & 3.67 & 0.998 & 0.032 bits \\
\bottomrule
\end{tabular}
\end{center}
Prime $3$ ($\alpha \approx 0.22$) is the bottleneck; primes $\ge 7$ decay as $(\log x)^{-3.5}$ or faster.

\subsection{Sieve diffusion model}
Starting from $f_0 = e^{-v}$ (the Cram\'er model~\cite{Cramer1936}) and iteratively adding sieve primes:
adding $p = 3$ alone reduces KL divergence from the empirical distribution by $74\%$ 
(from $0.061$ to $0.016$ nats at scale $10^7$).


% ====================================================================
\section{Discussion: What the Paper Does and Does Not Prove}
\label{sec:discussion}
% ====================================================================

\subsection{The role of the PDE model}

Gallagher~\cite{Gallagher1976} has shown (conditionally) that the distribution of prime gaps converges to Poisson. Our work does not 
provide another proof of this theorem; rather, it outlines a mechanism or an explanation for why this convergence occurs and how quickly it proceeds. 
As Tao~\cite{Tao2015} notes, probabilistic models used in number theory are ``largely informal and fuzzy in nature,'' yet 
they ``indirectly enable rigorous progress'' by serving as ``directional guides'' for making conjectures and as ``error-detection tools.'' 
Therefore, our use of probabilistic methods in number theory fits into a long line of research initiated by the random matrix theory program 
for $L$-functions. Keating and Snaith~\cite{KS2000}, using random matrix characteristic polynomials to model zeta functions, 
have predicted many properties of the zeta function and generated testable predictions for future verification. Similarly, we have derived the 
Burgers-Fisher equation from a particular form of a Fokker-Planck model, made predictions about its convergence rate (e.g., $\alpha \approx 0.19$), 
and verified those predictions computationally. However, whether our assumptions regarding the Fokker-Planck model can be rigorously derived 
from number theory remains an open question. The logical structure of our arguments is therefore:

\begin{enumerate}
\item \textbf{Pure PDE (Theorems~\ref{thm:burgers-fisher}--\ref{thm:convergence}):} A solution to the Feller Fokker-Planck equation represents the 
evolution of a probability density. Thus, if a probability density evolves under the action of a Feller Fokker-Planck operator, then there exists a 
velocity function satisfying the Burgers-Fisher equation. This velocity function has one stable equilibrium point located at $u=1$ and possesses a 
spectral gap equal to one. All nontrivial perturbations to the steady-state solution with decaying forcing will vanish over time. These are \emph{unconditional} theorems.
\item \textbf{Model (Hypothesis~\ref{hyp:sieve}):} The probability density of prime gaps near scale $s=\log x$ is accurately approximated by the 
stationary density associated with a Feller Fokker-Planck operator, perturbed by the Hardy-Littlewood singular series. Furthermore, the magnitude of 
these perturbations decreases as $O(1/s^{\beta})$, where $\beta > 0$. \emph{This is an empirical modeling assumption, which is supported by numerical 
computations but lacks theoretical derivation from first principles in number theory.}
\item \textbf{Consequences:} Using this model, we show that due to the existence of a spectral gap in the velocity function associated with the 
Burgers-Fisher equation, solutions to this equation corresponding to initial conditions given by prime-gap distributions evolve toward the Poisson distribution at a computable rate.
\end{enumerate}

Thus, although our results demonstrate the convergence of prime gaps to Poisson using a probabilistic framework modeled after fluid dynamics, 
establishing this result rigorously requires a derivation of Hypothesis~\ref{hyp:sieve} based upon the Hardy-Littlewood conjecture. 
This would represent an important step forward in closing the gap between number-theoretic models and their rigorous mathematical 
formulations---as was done by Talagrand in his derivation of the Parisi formula~\cite{Talagrand2006}, which had been predicted by physicists 27 years earlier.

\subsection{Open Problems}

We pose three open questions:
\begin{enumerate}
\item \textbf{Derive the Sieve Diffusion Hypothesis} from Hardy-Littlewood and the equidistribution of primes in arithmetic progressions.
\item \textbf{Show $\MI \to 0$} from the Burgers-Fisher framework (requires extending from marginal to joint gap densities).
\item \textbf{Compute the Exponent} $\alpha \approx 0.19$ from first principles, starting from pore strength $\delta_3 = 1$ and the equidistribution rate of gaps modulo $3$.
\end{enumerate}

% ====================================================================
\section*{Computational Details}
% ====================================================================

To run the numerical simulations, Python 3 is required, along with the SymPy, NumPy, and SciPy libraries. 
Code, data, and an interactive visualization can be found at \url{https://github.com/tobiascanavesi/prime-fluid-dynamics}. 
The interactive HTML was generated with the goal of bringing the ideas of the paper to a wider audience.

\subsection*{Use of Generative AI}
Claude Opus 4.6 (Anthropic, 2025) assisted in the development of code and formatting in \LaTeX{}. 
All mathematical content is the original contribution of the author.


\begin{thebibliography}{99}

\bibitem{Gallagher1976}
P.~X.~Gallagher, \emph{On the distribution of primes in short intervals}, Mathematika \textbf{23} (1976), 4--9.

\bibitem{HL1923}
G.~H.~Hardy and J.~E.~Littlewood, \emph{Some problems of `Partitio Numerorum': III}, Acta Math.\ \textbf{44} (1923), 1--70.

\bibitem{LO2016}
R.~J.~Lemke Oliver and K.~Soundararajan, \emph{Unexpected biases in the distribution of consecutive primes}, PNAS \textbf{113} (2016), E4446--E4454.

\bibitem{Holt2024}
F.~Holt, \emph{Expected biases in the distribution of consecutive primes}, preprint (2024), arXiv:2405.03540.

\bibitem{Szego1975}
G.~Szeg\H{o}, \emph{Orthogonal Polynomials}, 4th ed., AMS, 1975.

\bibitem{Risken1989}
H.~Risken, \emph{The Fokker--Planck Equation}, 2nd ed., Springer, 1989.

\bibitem{Cohen2024}
J.~E.~Cohen, \emph{Gaps between consecutive primes and the exponential distribution}, Experimental Math.\ \textbf{34}(2) (2025), 260--269.

\bibitem{KS2000}
J.~P.~Keating and N.~C.~Snaith, \emph{Random matrix theory and $\zeta(1/2+it)$}, Comm.\ Math.\ Phys.\ \textbf{214} (2000), 57--89.

\bibitem{Tao2015}
T.~Tao, \emph{254A, Supplement 4: Probabilistic models and heuristics for the primes}, blog post, What's New (2015). \url{https://terrytao.wordpress.com/2015/01/04/254a-supplement-4-probabilistic-models-and-heuristics-for-the-primes-optional/}

\bibitem{Talagrand2006}
M.~Talagrand, \emph{The Parisi formula}, Ann.\ of Math.\ \textbf{163} (2006), 221--263.

\bibitem{CIR1985}
J.~C.~Cox, J.~E.~Ingersoll, and S.~A.~Ross, \emph{A theory of the term structure of interest rates}, Econometrica \textbf{53}(2) (1985), 385--407.

\bibitem{Fisher1937}
R.~A.~Fisher, \emph{The wave of advance of advantageous genes}, Ann.\ Eugenics \textbf{7}(4) (1937), 355--369.

\bibitem{KPP1937}
A.~N.~Kolmogorov, I.~G.~Petrovsky, and N.~S.~Piskunov, \emph{\'Etude de l'\'equation de la diffusion avec croissance de la quantit\'e de mati\`ere et son application \`a un probl\`eme biologique}, Bull.\ Univ.\ Moscou, S\'er.\ Int., Sect.~A \textbf{1}(6) (1937), 1--25.

\bibitem{Cramer1936}
H.~Cram\'er, \emph{On the order of magnitude of the difference between consecutive prime numbers}, Acta Arith.\ \textbf{2} (1936), 23--46.

\end{thebibliography}

\end{document}
